\section{Linealización exacta por realimentación}

Se considera el modelo PVTOL
\begin{align}
  \ddot{x}      &= -\left(\frac{U}{m}\right)\sin(\theta), \label{eq:x}\\
  \ddot{z}      &=  \left(\frac{U}{m}\right)\cos(\theta)-g, \label{eq:z}\\
  \ddot{\theta} &= \frac{1}{J}\tau. \label{eq:theta}
\end{align}

Aquí, $\theta$ es el ángulo de actitud; $x$ y $z$ son las posiciones
traslacionales; $U$ es el empuje total; $\tau$ es el par aplicado para
controlar la actitud; $m$ es la masa; $J$ es el momento de inercia, y $g$ es
la aceleración de la gravedad. Se emplean
\[
  m=1,\qquad J=1.
\]
Por tanto,
\[
  \ddot{x}=-U\sin(\theta),\qquad
  \ddot{z}=U\cos(\theta)-g,\qquad
  \ddot{\theta}=\tau.
\]

Las aceleraciones traslacionales se escriben como
\[
  \begin{bmatrix}\ddot{x}\\\ddot{z}\end{bmatrix}
  =
  \begin{bmatrix}-\sin\theta\\\cos\theta\end{bmatrix}U+
  \begin{bmatrix}0\\-g\end{bmatrix}.
\]
Como ambas dependen de la misma entrada $U$, se definen las entradas virtuales
\[
  \boxed{v_x:=\ddot{x}},\qquad
  \boxed{v_z:=\ddot{z}}.
\]
El objetivo es imponer
\[
  \boxed{\ddot{x}=v_x,\qquad \ddot{z}=v_z}.
\]

Las entradas virtuales satisfacen
\[
  v_x=-U\sin\theta,\qquad
  v_z=U\cos\theta-g.
\]
Así,
\[
  v_x^2+(v_z+g)^2=U^2
\]
y, tomando $U\geq 0$,
\begin{equation}
  \boxed{U_d=\sqrt{v_x^2+(v_z+g)^2}}. \label{eq:Ud}
\end{equation}

Como las componentes del empuje requerido satisfacen:
\[
  -v_x=U_d\sin\theta_d,\qquad
  v_z+g=U_d\cos\theta_d.
\]

Por tanto,
\[
  \sin\theta_d=\frac{-v_x}{U_d},\qquad
  \cos\theta_d=\frac{v_z+g}{U_d},
\]
y el ángulo se obtiene mediante
\[
  \theta_d=
  \operatorname{atan2}\bigl(U_d\sin\theta_d,U_d\cos\theta_d\bigr)
\]
\[
  \boxed{\theta_d=\operatorname{atan2}(-v_x,v_z+g)}. \label{eq:theta-d}
\]